\section*{الفصل \ref{ch:b2:retrosynthesis} --- التخليق الراجع واستراتيجية الخطوات المتعددة}

\begin{solution}{exo:b2:retrosynthesis:1}
\begin{itemize}
  \item 2-فينيل بروبان-2-ول: البروبانون + \ce{PhMgBr}، أو الأسيتوفينون + \ce{CH3MgBr}.
  \item البنتان-3-ول: البروبانال + \ce{CH3CH2MgBr}.
  \item 1-هكسيل حلقي إيثانول: الإيثانال + بروميد هكسيل حلقي مغنيزيوم، أو هكسان حلقي كربالدهيد +
        \ce{CH3MgBr}.
  \item 2-ميثيل بيوتان-2-ول: البروبانون + \ce{CH3CH2MgBr}، أو البيوتانون + \ce{CH3MgBr}.
  \item 1-فينيل إيثانول: البنزالدهيد + \ce{CH3MgBr}، أو الإيثانال + \ce{PhMgBr}.
\end{itemize}
\end{solution}

\begin{solution}{exo:b2:retrosynthesis:2}
\omterm{def:b2:retrosynthesis:disconnection}{السنثون} المستقبل $\mathrm{Ph\overset{+}{C}H{-}OH}$، ومكافئه البنزالدهيد؛ \omterm{def:b2:retrosynthesis:disconnection}{والسنثون} المانح \ce{CH3-}،
ومكافئه بروميد ميثيل المغنيزيوم (أو ميثيل الليثيوم). (أما القطع الآخر، الذي يكون فيه \ce{Ph-} مانحًا، فيستعمل
\ce{PhMgBr} والإيثانال.)
\end{solution}

\begin{solution}{exo:b2:retrosynthesis:3}
$0.90 \times 0.85 \times 0.80 \times 0.95 \times 0.70 = 0.407$: أي نحو 41~\%.
\end{solution}

\begin{solution}{exo:b2:retrosynthesis:4}
\ce{LiAlH4}، اللازم لاختزال الإستر إلى الكحول، كان سيختزل الكيتون أيضًا. فاحمِ
الكيتون على هيئة أسيتاله الحلقي (الإيثان-1،2-ديول، وحفاز حمضي، مع إزالة الماء)؛ ثم اختزل الإستر بواسطة
\ce{LiAlH4}؛ ثم أزل الأسيتال بحمض مائي مخفف: 5-هيدروكسي بنتان-2-ون.
\end{solution}

\begin{solution}{exo:b2:retrosynthesis:5}
FGI: البيوتان-1،3-ديول $\Rightarrow$ 3-هيدروكسي بيوتانال (اختزال الألدهيد، \ce{NaBH4}). \omterm{def:b2:retrosynthesis:disconnection}{والفك} عند
مجموعتين لهذا الألدهيد $\beta$-الهيدروكسيلي عند الرابطة \babelsublr{$\alpha$--$\beta$}: جزيئان من الإيثانال
(\omterm{def:b2:enolates-aldol:aldol}{ألدول}). وإلى الأمام: الإيثانال، قاعدة مخففة، بارد؛ ثم \ce{NaBH4}.
\end{solution}

\begin{solution}{exo:b2:retrosynthesis:6}
تُقطع حلقة الهكسين الحلقي عند الرابطتين \ce{C-C} اللتين يكوّنهما \omterm{def:b2:frontier-orbitals:diels-alder}{تفاعل ديلز--ألدر}، أي
الأليليتين بالنسبة إلى الرابطة المزدوجة الحلقية في جهة البديل: البوتا-1،3-ديين والبوت-3-ين-2-ون، 
ومجموعة الأسيتيل على \omterm{def:b2:frontier-orbitals:diels-alder}{المحب للديين}.
\end{solution}

\begin{solution}{exo:b2:retrosynthesis:7}
بترقيم C1 (\ce{C=O})، وC2=C3، ثم C4 الحامل لمجموعتي الميثيل
(\cref{met:b2:conjugate-additions:robinson-disconnection}): القطع الألدولي C2--C3 يترك ألدهيدًا
عند C3 وميثيل كيتون؛ وقطع مايكل C4--C5 يترك 2-ميثيل بروبانال (C3 كربونيله، وC4 كربونه
$\alpha$) والبوت-3-ين-2-ون.
\end{solution}

\begin{solution}{exo:b2:retrosynthesis:8}
خطوة واحدة: البنزين مع كلوريد البيوتانويل و\ce{AlCl3} (\omterm{def:b2:aromatic-substitution:friedel-crafts}{أسيلة فريدل--كرافتس}، بلا إعادة ترتيب،
وأسيلة واحدة). وخطوتان: البنزالدهيد مع بروميد بروبيل المغنيزيوم، ثم أكسدة الكحول
الثانوي. وطريق فريدل--كرافتس أقصر وكواشفه أرخص، لكنه يستعمل مكافئًا كاملًا من
كلوريد الألومنيوم، الذي ينتهي في النفايات المائية.
\end{solution}

\begin{solution}{exo:b2:retrosynthesis:9}
كلوروكرومات البيريدينيوم في ثنائي كلوروميثان، أو أكسدة سويرن أو بيريودينان ديس--مارتن
تعطي الهكسانال؛ أما ثنائي الكرومات المائي المحمَّض أو البرمنغنات فتعطي حمض الهكسانويك، لأن
الألدهيد في الماء يكوّن هيدراته، الذي يتأكسد من جديد (\cref{prop:b2:retrosynthesis:mild-oxidation}).
\end{solution}

\begin{solution}{exo:b2:retrosynthesis:10}
يتفاعل الأمين أولًا، دون حماية، إذا كان الكحول محميًا: احمِ الكحول على هيئة إيثر سيليل
(يصمد أمام الأسيلة)؛ ثم أسِّل الأمين؛ ثم أزل مجموعة السيليل بالفلوريد؛ ثم أسِّل
الكحول. ولو وجب إبقاء الأمين حرًا بينما يتفاعل الكحول، لحُمي على هيئة كربامات
\textit{tert}-بيوتوكسي كربونيل (تُزال بالحمض)، المتعامدة مع إيثر السيليل (الذي يُزال
بالفلوريد): فيمكن تحرير أي منهما دون مسّ الآخر.
\end{solution}

\begin{solution}{exo:b2:retrosynthesis:11}
\ce{CH3COCH2CH2COCH3} مركب 1،4-ثنائي الكربونيل: اقطع الرابطة \ce{C-C} المركزية. المانح أيون إينولات
3-أوكسوبيوتانوات الإيثيل (إيثوكسيد الصوديوم)؛ والمستقبل، وهو كربون $\alpha$ لكيتون، يقدمه
1-بروموبروبان-2-ون، \ce{BrCH2COCH3} (استبدال ثنائي الجزيئية للبروميد). ثم تزيل الحلمهة المائية
للإستر والتسخين \ce{CO2}: الهكسان-2،5-ديون.
\end{solution}

\begin{solution}{exo:b2:retrosynthesis:12}
التخليق الراجع: الميثيل كيتون الذي بجواره \ce{CH2} هو ناتج تخليق
إستر الأسيتوأسيتيك؛ فاقطع الرابطة بين C3 وC4: أيون إينولات 3-أوكسوبيوتانوات الإيثيل والهاليد الأليلي
1-برومو-3-ميثيل بوت-2-ين، \ce{(CH3)2C=CHCH2Br}. وإلى الأمام: إيثوكسيد الصوديوم في الإيثانول، ثم الهاليد؛
ثم قاعدة مائية، وتحميض وتسخين، فتُحلمه الإستر ويُزال \ce{CO2}
(\cref{prop:b2:enolates-aldol:decarboxylation}): 6-ميثيل هبت-5-ين-2-ون.
\end{solution}

\begin{solution}{pb:b2:retrosynthesis:1}
\textbf{1.} \ce{HO-C6H4-CH2CH2COCH3}، ومجموعة \ce{OH} في الموضع بارا بالنسبة إلى السلسلة: فينول وكيتون.
\textbf{2.} هدرجة \ce{C=C}: \ce{ArCH2CH2CO} $\Rightarrow$ \ce{ArCH=CHCO}، وهو كيتون
$\alpha,\beta$-غير مشبع.
\textbf{3.} (\textit{E})-4-(4-هيدروكسي فينيل)بوت-3-ين-2-ون، \ce{HOC6H4CH=CHCOCH3}.
\textbf{4.} كربونيل $\alpha,\beta$-غير مشبع: \omterm{def:b2:enolates-aldol:aldol}{تكاثف ألدولي}، يُقطع عند \ce{C=C}.
\textbf{5.} 4-هيدروكسي بنزالدهيد والبروبانون.
\textbf{6.} ليس للألدهيد هيدروجين $\alpha$ فلا يكون إلا المحب للإلكترونات؛ والبروبانون،
بفائض، هو المصدر الوحيد \omterm{def:b2:enolates-aldol:enolate}{لأيون الإينولات} ويتفاعل مع الألدهيد الأكثر تفاعلية لا مع نفسه.
\textbf{7.} (1) 4-هيدروكسي بنزالدهيد، وفائض من البروبانون، وهيدروكسيد الصوديوم المائي؛ ثم التحميض. (2)
\ce{H2} على البلاديوم على الكربون، في درجة حرارة الغرفة، حتى يُمتص مكافئ واحد.
\textbf{8.} مجموعة \ce{CHO} تثبّت الفينوكسيد، فيكون هذا الفينول حمضيًا بقدر الفينول على الأقل
($\mathrm pK_a$ 9.99). وعند $\mathrm{pH} > 13$، يكون $[\ce{ArO-}]/[\ce{ArOH}] > 10^{13-10} = 10^3$: فهو
الفينوكسيد.
\textbf{9.} تدفع \ce{O-} الكثافة الإلكترونية عبر الحلقة نحو كربون الكربونيل (مانح
بالرنين في الموضع بارا بالنسبة إليه): فيصير الألدهيد أقل محبة للإلكترونات، والتكاثف أبطأ.
\textbf{10.} إيثر البنزيل: فالهدرجة على البلاديوم التي تختزل \ce{C=C} تشطر أيضًا
إيثر البنزيل (تحلل هيدروجيني)، فتحرر الفينول في الخطوة نفسها.
\textbf{11.} ثلاث: البنزلة، \omterm{def:b2:enolates-aldol:aldol}{والتكاثف الألدولي}، والهدرجة مع نزع الحماية.
\textbf{12.} لا: ففي الظروف اللطيفة (درجة حرارة الغرفة، البلاديوم) لا تُختزل الحلقة \omterm{def:b2:huckel:aromatic}{العطرية}.
\textbf{13.} عدم الحماية: فالفينوكسيد لا يفعل سوى إبطاء التكاثف، وهو ما يعوّضه زمن أطول أو
تسخين؛ وتوفير خطوة يساوي أكثر من الربح في السرعة.
\textbf{14.} 4-يودوفينول والبوت-3-ين-2-ون.
\textbf{15.} \ce{HOC6H4I + CH2=CHCOCH3 + Et3N -> HOC6H4CH=CHCOCH3 + Et3NH+ + I-}، بحفاز من البلاديوم.
\textbf{16.} عند \ce{CH2} الطرفية، أي الكربون $\beta$: فمجموعة الأريل تدخل عند الطرف الأقل
إعاقة، ثم يعطي حذف الهيدريد $\beta$ الإينون (\textit{E}) المترافق.
\textbf{17.} \omterm{def:b2:catalytic-cycles:oxidative-addition}{الإضافة التأكسدية} ليوديد الأريل إلى \ce{Pd(0)}، وتناسق الألكين وإدخاله،
وحذف الهيدريد $\beta$، وتجديد \ce{Pd(0)} بالقاعدة، التي تأخذ \ce{HI}
(\cref{ch:b2:catalytic-cycles}).
\textbf{18.} $162.188/(220.005 + 70.091 + 101.193) = 162.188/391.289 \approx 41.4~\%$.
\textbf{19.} \ce{C7H6O2 + C3H6O -> C10H10O2 + H2O}: $162.188/(122.123 + 58.080) = 162.188/180.203
\approx 90.0~\%$.
\textbf{20.} 100~\%: فكل ذرة من \ce{H2} ومن الإينون تنتهي في الناتج.
\textbf{21.} $164.204/(122.123 + 58.080 + 2.016) = 164.204/182.219 \approx 90.1~\%$.
\textbf{22.} الرابطة \ce{C=C}: فهي مترافقة وغير معاقة، فترتبط وتُهدرج على البلاديوم أسرع
بكثير من \ce{C=O} الكيتون؛ والتوقف بعد مكافئ واحد من الهيدروجين يحفظ الكيتون.
\textbf{23.} الطريق الألدولي: ألدهيد رخيص، والبروبانون وهيدروكسيد الصوديوم، والماء ناتج ثانوي. وطريق
هيك: يوديد أريل، وحفاز من البلاديوم، والبوت-3-ين-2-ون، الشديد السمية.
\textbf{24.} الطريق الألدولي: خطوتان، وكواشف أرخص وأكثر أمانًا، واقتصاد ذري فوق 90~\%.
\textbf{25.} تضربه في 0.90: $0.782 \times 0.90 \approx 70~\%$.
\textbf{26.} $0.85 \times 0.92 = \boldsymbol{\approx 78~\%}$ إجماليًا.
\end{solution}
